\documentclass[10pt,a4paper]{amsart}
\usepackage[margin=19mm]{geometry}
\usepackage{amsmath,amssymb,mathtools}
\usepackage[scr=rsfs,cal=euler]{mathalfa}
\usepackage{xfrac}
\usepackage[unicode,hidelinks,bookmarks=false]{hyperref}
\newcommand{\ie}{i.e.\ }
\newcommand{\cf}{cf.\ }
\newcommand{\resp}{respectively\ }
\newcommand{\Subsection}{\subsection}
\providecommand{\nobreakdash}{\nobreak\hbox{-}}
\setlength{\emergencystretch}{2em}
\multlinegap0pt
\usepackage{amssymb}
\usepackage{xfrac}
\usepackage{float}
\newtheorem{corollary}{Corollary}
\newtheorem{theorem}{Theorem}
\newtheorem{remark}{Remark}
\newtheorem{lemma}{Lemma}
\newcommand{\F}{\mathbb{F}}
\newcommand{\Q}{\mathbb{Q}}
\newcommand{\Z}{\mathbb{Z}} 
\newcommand{\rank}{\mathrm{rank}}
\title{Algorithms for $p$-rationality\\ and for $p$-saturation of units}
\author{Tommy Hofmann, Henri Johnston}
\date{Source: arXiv:2609.15190v1 (CC BY 4.0)}
\begin{document}
\maketitle

\noindent\emph{Source and licence.} Algorithms for $p$-rationality\\ and for $p$-saturation of units. Version arXiv:2609.15190v1 (CC BY 4.0), distributed by arXiv under the Creative Commons Attribution 4.0 International licence, http://creativecommons.org/licenses/by/4.0/, as recorded in the arXiv licence metadata for this exact version and shown on the abstract page https://arxiv.org/abs/2609.15190v1. Full licence notice: \texttt{licenses/arxiv-2609.15190-CC-BY-4.0.txt}.\par\smallskip

\noindent\textit{Editorial scope.} Counted unit: the article's theorem thm:p-odd-rat-cyclo (source0.tex lines 960--967). The article's complete original proof of it is delivered with it (source0.tex lines 969--1039, one \texttt{proof} environment of 71 lines). No mathematics is written, repaired or re-derived: the statement and the proof are the article's own lines, in the article's own notation, and no retained line is substituted or re-flowed.  Retained uncounted as the source-local context the proof needs: lines 500--524; lines 690--701; lines 711--734; lines 919--929; lines 934--937.  Nothing was omitted from any retained span, so the package carries no omission record.  No retained line contains a citation, so the wrapper carries no bibliography and no bibliography entry.  No retained line contains a figure environment, an image-inclusion command or drawing code, so no image is delivered and the package declares no asset dependency.  Editorial formatting only, in addition: an AMS \texttt{amsart} preamble that restates the article's own declarations and macros used by the retained lines, namely the environments \texttt{corollary}, \texttt{theorem}, \texttt{remark}, \texttt{lemma} on separate counters, together with the standard packages those lines need.  The retained environments print in this document as Corollary 1, Theorem 1, Remark 1, Lemma 1 in order of appearance, whereas the article prints the same items with the numbering of its own full document.  The complete native source of the article as distributed in the arXiv e-print for this exact version is bundled unchanged as \texttt{sources/210395/source0.tex}.
\begin{corollary}\label{cor:p-sat}
Let $B$ be a subgroup of a finitely generated abelian group $A$.  
Let $p$ be a prime number and assume that $A[p] \leq B$.
Let $X$ be an $\F_p$-vector space and let $\psi \colon A \to X$ be a group homomorphism. 
Let $\psi \vert_{B}$ be the restriction of $\psi$ to $B$.
\renewcommand{\labelenumii}{(\roman{enumii})}
\begin{enumerate}
    \item 
    We have $B^{p} \leq \ker(\psi \vert_{B})$, and hence $\psi \vert_{B}$
    induces an $\F_{p}$-linear map $\overline{\psi \vert_{B}} \colon B/B^{p} \to X$.
    \item The following are equivalent:
    \begin{enumerate}
        \item $B^{p} = \ker(\psi \vert_{B})$,
        \item $\dim_{\F_p}(\psi(B)) = \rank_{p}(B)$,
        \item $\dim_{\F_p}(\overline{\psi \vert_{B}}(B/B^{p})) = \rank_{p}(B)$,
        \item $\overline{\psi \vert_{B}}$ is injective.
    \end{enumerate}
    \item If the equivalent conditions of \textup{(b)} 
    hold then $B$ is $p$-saturated in $A$.
    \item If $\dim_{\F_p}(\psi(B)) = \rank_{p}(A)$ then $B$ is $p$-saturated and of finite index in $A$.
    \item If $b \in B$ and $b = a^{p}$ for some $a \in A \setminus B$, 
    then $bB^{p}$ is a non-trivial element of $\ker(\overline{\psi \vert_{B}})$, and 
    $[\langle B, a \rangle : B]=p$.
\end{enumerate}    
\end{corollary}

\begin{theorem}\label{thm:nqd-char-p-rat}
A number field $K$ is $p$-rational precisely when all three of the following conditions are
satisfied:
\begin{enumerate}
    \item The compositum of all $\Z_{p}$-extensions of $K$ contains $H_{p}(K)$.
    \item The diagonal map $\lambda_{K,p} : \mu_{p}(K) \longrightarrow \prod_{\mathfrak{p} \in S_{p}(K)} 
    \mu_{p}(K_{\mathfrak{p}})$ is an isomorphism.
    \item The diagonal map 
    $\overline{\eta}_{K,p} : \mathcal{O}_{K}^{\times}/\mathcal{O}_{K}^{\times p} \longrightarrow 
    \prod_{\mathfrak{p} \in S_{p}(K)} U_{\mathfrak{p}}/ U_{\mathfrak{p}}^{p}$ is injective.
\end{enumerate}
\end{theorem}

\begin{remark}\label{rmk:roots-of-unity-condition}
Condition (b) is equivalent to the following pair of assertions:
\[
\begin{cases}
\lvert S_{p}(K) \rvert =1, & \text{ if } \mu_{p} \subset K,  \\
\mu_{p} \not \subset K_{\mathfrak{p}} \text{ for every } \mathfrak{p} \in S_{p}(K),
& \text{ if } \mu_{p} \not \subset K.
\end{cases}
\]
For $\mathfrak{p} \in S_{p}(K)$, let $e_{\mathfrak{p}}$ denote the
ramification index of $\mathfrak{p}$ in $K/\Q$.
If $(p-1) \nmid e_{\mathfrak{p}}$ then $\mu_{p} \not \subset K_{\mathfrak{p}}$ 
(and so $\mu_{p} \not \subset K$) because 
$\Q_{p}(\zeta_{p})/\Q_{p}$ is totally ramified and of degree $p-1$.
In particular, if $p>[K:\Q]+1$ or $p \nmid 2d_{K}$ then condition (b) holds.

If $\mu_{p} \not \subset K$ then condition (b) is equivalent to the following:
\begin{itemize}
    \item no prime $\mathfrak{p} \in S_{p}(K)$ splits completely in $K(\zeta_{p})/K$.
\end{itemize}
In fact, by the above discussion, it is only necessary to check this
condition for $\mathfrak{p} \in S_{p}(K)$ such that $(p-1) \mid e_{\mathfrak{p}}$.
This observation is useful for computational purposes.
\end{remark}

\begin{theorem}\label{thm:sinnott-cyclo}
Let $n \in \Z_{\geq 3}$ with $n \not \equiv 2 \bmod 4$ and
let $g$ be the number of distinct prime factors of $n$.
Let $b=0$ if $g=1$ and $b=2^{g-2}+1-g$ if $g \geq 2$.
Then
\[
[E_{n} : C_{n}] 
=[E_{n}^{+} : C_{n}^{+}] 
= 2^{b}h_{n}^{+}.
\]
\end{theorem}

\begin{lemma}\label{lem:real-cyclo-p-rat-class-number}
Let $n \in \Z_{\geq 3}$ with $n \not \equiv 2 \bmod 4$.
If $\Q(\zeta_{n})^{+}$ is $p$-rational then $p \nmid h_{n}^{+}$.
\end{lemma}

\begin{theorem}\label{thm:p-odd-rat-cyclo}
Let $n \in \Z_{\geq 3}$ with $n \not \equiv 2 \bmod 4$ and let $p$ be an odd prime number.
Then $\Q(\zeta_{n})^{+}$ is $p$-rational if and only if
\begin{enumerate}
    \item either $p \nmid n$ or no prime of $\Q(\zeta_{n})^{+}$ above $p$ splits in $\Q(\zeta_{n})/\Q(\zeta_{n})^{+}$, and
    \item $\dim_{\F_{p}}(\eta_{\Q(\zeta_{n})^{+},p}(C_{n}^{+}))=\frac{1}{2}\varphi(n)-1$.
\end{enumerate}
\end{theorem}

\begin{proof}
Henceforth abbreviate $\eta_{\Q(\zeta_{n})^{+},p}$, $\overline{\eta}_{\Q(\zeta_{n})^{+},p}$
and $\lambda_{\Q(\zeta_{n})^{+},p}$ to 
$\eta$, $\overline{\eta}$ and $\lambda$,
respectively.
Since $p$ is odd and $\Q(\zeta_{n})^{+}$ is totally real, we have
\begin{equation}\label{eq:p-odd-unit-rank-real-cyclo}
\rank_{p}(E_{n}^{+}) 
= \rank_{\Z}(E_{n}^{+})
= [\Q(\zeta_{n})^{+}:\Q] -1 
= \textstyle{\frac{1}{2}}\varphi(n)-1.
\end{equation}
If $p \mid n$ then $\Q(\zeta_{n})^{+}(\zeta_{p})=\Q(\zeta_{n})$ and if $p \nmid n$ then
$p \nmid d_{\Q(\zeta_{n})^{+}}$. Hence by Remark~\ref{rmk:roots-of-unity-condition}, 
condition 
(a) is equivalent to the bijectivity of $\lambda$.

Suppose that $\Q(\zeta_{n})^{+}$ is $p$-rational. 
By Theorem~\ref{thm:nqd-char-p-rat} 
the map $\lambda$ is an isomorphism and so condition (a) holds
by the discussion above.
Lemma~\ref{lem:real-cyclo-p-rat-class-number} implies that $p \nmid h_{n}^{+}$.
Thus $p \nmid [E_{n}^{+} : C_{n}^{+}]$ by Theorem~\ref{thm:sinnott-cyclo} and
so $\eta(C_{n}^{+}) = \eta(E_{n}^{+})$.
Again by Theorem~\ref{thm:nqd-char-p-rat}, $\overline{\eta}$ is injective
and thus by the rank-nullity theorem, the previous sentence and \eqref{eq:p-odd-unit-rank-real-cyclo} we have 
\[
\dim_{\F_{p}}(\eta(C_{n}^{+}))
=
\dim_{\F_{p}}(\eta(E_{n}^{+}))
=
\dim_{\F_{p}}(\overline{\eta}(E_{n}^{+}/(E_{n}^{+})^{p}))
=
\rank_{p}(E_{n}^{+})
= 
\textstyle{\frac{1}{2}}\varphi(n)-1.
\]
Therefore condition (b) also holds.

Suppose conversely that (a) and (b) hold. 
Then $\lambda$ is an isomorphism by (a) and the discussion above.
By (b) and \eqref{eq:p-odd-unit-rank-real-cyclo} we have
\[
\rank_{p}(E_{n}^{+})
=
\textstyle{\frac{1}{2}}\varphi(n)-1
= 
\dim_{\F_{p}}(\eta(C_{n}^{+}))
\leq 
\dim_{\F_{p}}(\eta(E_{n}^{+})) 
\leq 
\rank_{p}(E_{n}^{+}),
\]
and so these inequalities are in fact equalities. 
In particular, 
\[
\dim_{\F_{p}}(\overline{\eta}(E_{n}^{+}/(E_{n}^{+})^{p})) 
= \dim_{\F_{p}}(\eta(E_{n}^{+}))
=\rank_{p}(E_{n}^{+}),
\]
and so $\overline{\eta}$ is injective.
Moreover, 
$\dim_{\F_{p}}(\eta(C_{n}^{+}))=\rank_{p}(E_{n}^{+})$
and $E_{n}^{+}$ has no non-trivial $p$-torsion since $p$ is odd and $\Q(\zeta_{n})^{+}$
is totally real. Hence by Corollary~\ref{cor:p-sat}(d), 
$C_{n}^{+}$ is $p$-saturated and of finite index in $E_{n}^{+}$, 
or equivalently, $p \nmid [E_{n}^{+}:C_{n}^{+}]$.
Thus $p \nmid h_{n}^{+}$ by Theorem~\ref{thm:sinnott-cyclo},
and so $H_{p}(\Q(\zeta_{n})^{+})=\Q(\zeta_{n})^{+}$.
Therefore $\Q(\zeta_{n})^{+}$ is $p$-rational by Theorem~\ref{thm:nqd-char-p-rat}.
\end{proof}

\end{document}
